\section{Capítulo~\ref{sec:serocomputer}. Neurônios \nt{SERO}}

As propriedades dos próprios neurônios \nt{SERO} (ritmo, autoinibição, sinal definido pelo receptor, subsistemas) foram medidas diretamente; os cálculos do capítulo (regulador, filtro, lógica) decorrem de suas equações; $\tau_c$, o coeficiente de difusão da serotonina e o número de locais de liberação são estimativas; o papel do laço como regulador de nível no cérebro é uma hipótese (no núcleo da rafe a autoinibição é pontual e dá só pausas curtas), e o papel do \nt{SERO} como horizonte é uma hipótese com experimentos comportamentais a seu favor.

{\footnotesize
\setlength{\tabcolsep}{3pt}\renewcommand{\arraystretch}{1.15}
\begin{longtable}{>{\raggedright}p{0.5cm}>{\raggedright}p{4.0cm}>{\raggedright}p{5.6cm}>{\raggedright}p{2.3cm}>{\raggedright\arraybackslash}p{1.7cm}}
\toprule
N.º & Proposição & Experimento e o que foi medido & Fonte & Status \\
\midrule\endfirsthead
\toprule
N.º & Proposição & Experimento e o que foi medido & Fonte & Status \\
\midrule\endhead
1 & O sinal da ação da serotonina é escolhido pelo receptor do destinatário (proposição~\ref{prop:receptor}) & Os receptores de serotonina se dividem em famílias por farmacologia e mecanismo: o 5-HT$_{1A}$, via proteína G, abre canais de potássio e inibe a célula; o 5-HT$_{2A}$ e o ionotrópico 5-HT$_3$ excitam. Um mesmo neurotransmissor produz respostas opostas em células diferentes. & \cite{BarnesSharp1999} & medido \\
2 & Na vigília, o neurônio \nt{SERO} emite regularmente alguns disparos por segundo & Registro de neurônios do núcleo dorsal da rafe em gatos em livre movimento: descarga rítmica lenta de $2.8\pm0.2$ disparos/s na vigília tranquila; no sono de ondas lentas a taxa cai $34$--$68\,\%$, no sono REM, $98\,\%$. Em ratos anestesiados, $1.7\pm0.5$ disparos/s, muito regular. & \cite{TrulsonJacobs1979, GagnonParent2014, JacobsAzmitia1992} & medido \\
3 & O neurônio \nt{SERO} é um marca-passo: não precisa de empurrão a cada disparo, mas precisa de um aporte de fundo constante (na fatia, noradrenalina via receptores $\alpha_1$; no organismo, do \nt{NORA}) & Na fatia de cérebro, as células descarregam de forma lenta e regular; após cada disparo há uma pós-hiperpolarização de $200$--$800\ms$. Na fatia há menos células com atividade espontânea que no organismo: o ritmo regular precisa de uma entrada tônica de noradrenalina via receptores $\alpha_1$, cujo bloqueio o apaga. & \cite{VandermaelenAghajanian1983} & medido \\
4 & Quase todos os receptores são metabotrópicos; a resposta é lenta: sobe e desce em cerca de um segundo & Todas as famílias de receptores, exceto a 5-HT$_3$, são acopladas a proteínas G. Na fatia, a liberação evocada de serotonina produz, via 5-HT$_{1A}$, uma corrente inibitória que sobe e desce em um segundo. & \cite{BarnesSharp1999, CourtneyFord2016} & medido \\
5 & Nas regiões de destino, a serotonina sai no volume, e não só na fenda; no próprio núcleo da rafe, a inibição dos vizinhos é ponto a ponto & Voltametria cíclica rápida em fatias: a liberação por disparo é igual numa região com sinapses e no núcleo da rafe, onde quase não há sinapses; não depende do bloqueio da recaptação e dos receptores; há muito menos moléculas por terminal que transportadores e receptores, e a concentração extracelular coincide com a afinidade do receptor principal. No núcleo da rafe, a inibição via 5-HT$_{1A}$ é ponto a ponto: o transportador não deixa a serotonina acumular fora da fenda. & \cite{Bunin1998, CourtneyFord2016} & medido \\
6 & A concentração em~\eqref{eq:seroneuron} diminui pela recaptação; $\tau_c$ da ordem de um segundo é uma estimativa & Voltametria no cérebro vivo: a serotonina extracelular é limitada sobretudo pela recaptação e pela degradação, não pela síntese. Não há medição direta de $\tau_c$ nos trabalhos citados; a ordem de grandeza é estimativa do capítulo. & \cite{Hashemi2012serotonin} & medido; $\tau_c$: estimativa \\
7 & NÃO e NOU num único marca-passo; completude da lógica \nt{SERO} (proposição~\ref{prop:serocomplete}, fig.~\ref{fig:serogates}) & Decorre de~\eqref{eq:seroneuron}: um marca-passo inibível é um NOU, e o NOU é funcionalmente completo. Não há experimento em que se tenha montado lógica com neurônios \nt{SERO}; a condição de volumes separados não se cumpre no cérebro. & dedução no capítulo & teoria \\
8 & A serotonina inibe o próprio neurônio \nt{SERO} por receptores nele & Em ratos anestesiados, agonistas de 5-HT$_{1A}$ por via intravenosa e iontoforética inibem a descarga dos neurônios da rafe dorsal com a mesma relação dose-resposta da própria serotonina; agonistas de 5-HT$_{1B}$ quase não agem. Na fatia, a serotonina endógena das células vizinhas produz, via 5-HT$_{1A}$, uma corrente e uma pausa curta na descarga, sem mudar a taxa média. & \cite{Sprouse1987, CourtneyFord2016} & medido \\
9 & No modelo, o laço por um volume comum é um regulador de nível: a dispersão e a constante de tempo caem $1+G$ vezes~\eqref{eq:feedback}; que o laço funcione assim no cérebro é hipótese & Fórmula clássica do amplificador com realimentação negativa; deduzida de novo no capítulo. O ganho de laço $G$ do sistema \nt{SERO} não foi medido. Medido: na fatia, a autoinibição é ponto a ponto, dá uma pausa curta e não muda a taxa média. & \cite{Black1934feedback, CourtneyFord2016} & teoria; no cérebro, hipótese \\
10 & A concentração é uma média móvel da taxa, um filtro passa-baixas~\eqref{eq:lowpass}, fig.~\ref{fig:lowpass} & Solução da equação linear~\eqref{eq:seroneuron}; a figura é um cálculo por essa fórmula com $\tau_c=1\,\text{s}$. Não há modelo separado em \texttt{models/}. & dedução no capítulo & teoria \\
11 & A tortuosidade dos espaços torna a difusão de uma molécula pequena cerca de $2.6$ vezes mais lenta & Método da fonte pontual: iontoforese de tetrametilamônio e registro de sua concentração com eletrodo íon-seletivo em fatias e no cérebro vivo. Tortuosidade $\lambda\approx1.6$, ou seja, o $D$ efetivo é $\lambda^2\approx2.6$ vezes menor que o livre; fração de volume extracelular de cerca de $0.2$. O coeficiente de difusão da própria serotonina no tecido não foi medido nesses trabalhos; no capítulo, é uma estimativa. & \cite{Sykova2008} & medido \\
12 & Comprimento de um ``fio'' independente $\ell\sim\sqrt{D\tau}\approx20$~\textmu m~\eqref{eq:diffusionlength}; o número de fios é limitado pelo número dessas regiões (proposição~\ref{prop:volumewires}) & Lei da difusão e substituição das estimativas de $D$ e $\tau$ (linhas~6 e~11); a proposição~\ref{prop:volumewires} é uma consequência. Não há experimento direto que conte os canais independentes da transmissão por volume. & dedução no capítulo & teoria \\
13 & A saída de um neurônio \nt{SERO} se espalha por enormes áreas do cérebro; locais de liberação, por estimativa, $\sim10^5$ & Reconstrução completa de neurônios isolados da rafe dorsal no rato: todos os axônios ascendentes se ramificam muito, o comprimento total do axônio de uma célula chega a $18.7$~cm, e os ramos de uma única célula alcançam o estriado, o córtex pré-frontal e a amígdala. O número de locais de liberação por célula não foi contado diretamente. & \cite{GagnonParent2014} & medido; $10^5$: estimativa \\
14 & Os neurônios \nt{SERO} se dividem em alguns subsistemas com saídas e respostas diferentes & Marcação viral-genética em camundongos: os neurônios que vão para a amígdala e para o córtex frontal quase não se sobrepõem na ramificação, recebem entradas diferentes e respondem de modo oposto a um estímulo aversivo; ativar e desativar cada grupo muda o comportamento de modos diferentes. & \cite{Ren2018} & medido \\
15 & Condição de paciência $D<\tau_\gamma\ln(R/r_0)$~\eqref{eq:patience} com desconto exponencial; com o desconto medido, hiperbólico, $D<\tau_\gamma(R/r_0-1)$ & Decorre do desconto $e^{-D/\tau_\gamma}$ ou $1/(1+D/\tau_\gamma)$; dedução no capítulo. Macacos, escolha com atrasos de segundos: o desconto fica mais perto da hipérbole que da exponencial; a resposta dos neurônios \nt{DOPA} ao sinal antecipador de uma recompensa adiada também cai com o atraso de forma hiperbólica. & dedução no capítulo; \cite{KobayashiSchultz2008} & teoria \\
16 & O \nt{SERO} define o horizonte $\tau_\gamma$ (seção~\ref{sec:horizon}) & A ativação optogenética dos neurônios \nt{SERO} da rafe dorsal em camundongos prolonga a espera por uma recompensa adiada e aumenta a persistência na busca de uma recompensa que ficou rara. Em humanos, a redução da serotonina (dieta sem triptofano) torna mais íngreme o desconto da recompensa adiada. Como a concentração se converte em $\tau_\gamma$ não se sabe. & \cite{Doya2002, Miyazaki2014, Lottem2018, Schweighofer2008} & hipótese \\
\bottomrule
\end{longtable}}
